\unit{Exposé 190 -- Technique de descente I}
\sid{190}
\src{299}
\seminar{12\textsuperscript{e} année, 1959/60, n\textsuperscript{o} 190\\Décembre 1959}
\paperhead{TECHNIQUE DE DESCENTE ET THÉORÈMES D'EXISTENCE EN GÉOMÉTRIE ALGÉBRIQUE\\
I. GÉNÉRALITÉS. DESCENTE PAR MORPHISMES FIDÈLEMENT PLATS}%
  {par Alexander GROTHENDIECK}

D'un point de vue technique, le présent exposé, et ceux qui lui feront suite,
peuvent être considérés comme des variations sur le célèbre « théorème 90 » de
Hilbert. L'introduction de la méthode de descente en géométrie algébrique semble
due à A. Weil, sous le nom de « descente du corps de base ». Weil se restreint
d'ailleurs au cas d'extensions finies séparables de corps. Le cas d'extensions
radicielles de hauteur 1 a été ensuite traité par P. Cartier. Faute du langage
des schémas, et plus particulièrement faute d'éléments nilpotents dans les
anneaux qu'on se permettait de considérer, l'identité essentielle entre ces
deux cas n'avait pu être formulée par Cartier.

À l'heure actuelle, il semble que la technique générale de descente qui sera
exposée (jointe le cas échéant aux théorèmes fondamentaux de la « géométrie
formelle », cf. [3]) est à la base de la plupart des théorèmes d'existence en
géométrie algébrique.\footnote{L'erratum de 1962 juge cette formulation
excessive : elle vaut dans une large mesure pour les techniques non projectives
des deux premiers exposés de la série TDTE I à VI, mais non pour les techniques
projectives (TDTE IV, V et VI).} Il convient de noter d'ailleurs que ladite
technique peut certainement se transposer en « géométrie analytique », et on
peut espérer que dans un proche avenir, les spécialistes sauront démontrer les
analogues « analytiques » des théorèmes d'existence en géométrie formelle qui
seront donnés dans l'exposé II. En tout cas, les récents travaux de
Kodaira--Spencer, dont les méthodes semblent inaptes à la définition et l'étude
de « variétés de modules » au voisinage de leurs points singuliers, montrent
assez clairement la nécessité de méthodes plus proches de la théorie des
schémas (que devront compléter bien entendu les techniques transcendantes).

Dans le présent exposé I, nous traiterons le cas de descente le plus élémentaire
(indiqué dans le titre). Les applications des théorèmes 1, 2, 3 ci-dessous sont
cependant déjà fort nombreuses. Nous nous sommes bornés à en donner quelques-unes
à titre d'exemple, sans essayer de leur donner ici le maximum de généralité
possible.

Nous utiliserons librement le langage des schémas, pour lequel nous renvoyons à
l'exposé déjà cité, ainsi qu'à [2]. Notons cependant expressément que les
préschémas envisagés dans le présent exposé ne sont pas nécessairement
noethériens, et les

\src{300}
morphismes ne sont pas nécessairement de type fini. Ainsi, si $A$ est un anneau
local noethérien, de complété $\widehat A$, on aura à considérer l'anneau non
noethérien $\widehat A\otimes_A\widehat A$, et les morphismes de schémas affines
correspondants aux inclusions entre les anneaux envisagés.

\fsection{A}{Préliminaires sur les catégories}

\subsection*{1. Catégories fibrées, données de descente, morphismes de
$\mathcal F$-descente}
\addcontentsline{toc}{subsection}{1. Catégories fibrées et données de descente}

\result{a. DÉFINITION}{1.1} Une \emph{catégorie fibrée} $\mathcal F$ de base
$\mathcal C$ est formée d'une catégorie $\mathcal C$, pour tout
$X\in\mathcal C$ d'une catégorie $\mathcal F_X$, pour tout
$\mathcal C$-morphisme $f:X\longrightarrow Y$ d'un foncteur
$f^*:\mathcal F_Y\longrightarrow\mathcal F_X$, noté aussi
\[
  f^*(\xi)=\xi\times_Y X
\]
pour $\xi\in\mathcal F_Y$ ($X$ étant considéré comme un « objet de
$\mathcal C$ au-dessus de $Y$ », i.e. comme muni du morphisme $f$), et enfin,
pour deux morphismes composables
$X\xrightarrow{f}Y\xrightarrow{g}Z$, d'un isomorphisme de foncteurs
\[
  c_{f,g}:(gf)^*\longrightarrow f^*g^*,
\]
ces données étant assujetties aux conditions suivantes :

\begin{enumerate}
\item[(i)] $\operatorname{id}_X^*=\operatorname{id}$ ;
\item[(ii)] $c_{f,g}$ est l'isomorphisme identique si $f$ ou $g$ est un
isomorphisme identique ;
\item[(iii)] étant donnés trois morphismes composables
$X\xrightarrow{f}Y\xrightarrow{g}Z\xrightarrow{h}T$, on a commutativité dans
le diagramme suivant d'isomorphismes formé au moyen des isomorphismes de la
forme $c_{u,v}$ :
\[
\begin{array}{ccc}
  (h(gf))^*=((hg)f)^* & \longrightarrow & f^*(hg)^* \\
  \big\downarrow && \big\downarrow \\
  (gf)^*h^* & \longrightarrow & (f^*g^*)h^*=f^*(g^*h^*).
\end{array}
\]
\end{enumerate}

\result{EXEMPLE}{1} Soit $\mathcal C$ une catégorie où les produits fibrés
existent. On définit alors une catégorie fibrée $\mathcal F$ de base
$\mathcal C$ en posant $\mathcal F_X=$ catégorie des objets de $\mathcal C$
au-dessus de $X$, le foncteur
$f^*:\mathcal F_Y\longrightarrow\mathcal F_X$ correspondant à un morphisme
$f:X\longrightarrow Y$ étant défini par le produit fibré
$Z\longmapsto Z\times_YX$.

\result{EXEMPLE}{2} Soit $\mathcal C$ la catégorie des préschémas, et pour
$X\in\mathcal C$, soit $\mathcal F_X$ la catégorie des faisceaux
quasi-cohérents de modules sur $X$. Si $f:X\longrightarrow Y$ est un morphisme
de préschémas, $f^*:\mathcal F_Y\longrightarrow\mathcal F_X$ est le foncteur
image inverse

\src{301}
de faisceaux de modules. On obtient ainsi une catégorie fibrée de base
$\mathcal C$.

\result{b. DÉFINITION}{1.2} Un diagramme
\[
  E\xrightarrow{u}E'
  \overset{v_1}{\underset{v_2}{\rightrightarrows}}E''
\]
d'applications d'ensembles est dit \emph{exact} si $u$ est une bijection de
$E$ sur la partie de $E'$ formée des $x'\in E'$ tels que
$v_1(x')=v_2(x')$.

\result{DÉFINITION}{1.3} Soit $\mathcal F$ une catégorie fibrée de base
$\mathcal C$ ; considérons un diagramme
\[
  \tag{+}
  S\xleftarrow{\alpha}S'
  \overset{\beta_1}{\underset{\beta_2}{\rightrightarrows}}S''
\]
de morphismes dans $\mathcal C$ tels que
$\alpha\beta_1=\alpha\beta_2$ ; ce diagramme est dit
$\mathcal F$-\emph{exact} si, pour tout couple $\xi,\eta$ d'éléments de
$\mathcal F_S$, le diagramme suivant d'applications d'ensembles
\[
  \tag{+'}
  \operatorname{Hom}(\xi,\eta)
  \xrightarrow{\alpha^*}
  \operatorname{Hom}(\alpha^*\xi,\alpha^*\eta)
  \overset{\beta_1^*}{\underset{\beta_2^*}{\rightrightarrows}}
  \operatorname{Hom}(\gamma^*\xi,\gamma^*\eta),
\]
où $\gamma=\alpha\beta_1=\alpha\beta_2$, est exact. Dans ce dernier
diagramme, nous avons identifié $\beta_i^*\alpha^*$ à
$(\alpha\beta_i)^*=\gamma^*$ grâce à $c_{\beta_i,\alpha}$, pour simplifier.

\result{DÉFINITION}{1.4} Soit $\mathcal F$ une catégorie fibrée de base
$\mathcal C$ ; considérons deux morphismes $\beta_1,\beta_2$ de $S''$ dans
$S'$, dans $\mathcal C$. Soit
$\xi'\in\mathcal F_{S'}$. On appelle \emph{donnée de recollement} sur $\xi'$
(relativement au couple $(\beta_1,\beta_2)$) un isomorphisme de
$\beta_1^*(\xi')$ sur $\beta_2^*(\xi')$. Si $\xi',\eta'\in\mathcal F_{S'}$
sont munis chacun d'une donnée de recollement, un morphisme
$u:\xi'\longrightarrow\eta'$ dans $\mathcal F_{S'}$ est dit
\emph{compatible avec les données de recollement} si on a commutativité dans
le diagramme suivant :
\[
\begin{array}{ccc}
  \beta_1^*(\xi') & \longrightarrow & \beta_2^*(\xi') \\
  \big\downarrow && \big\downarrow \\
  \beta_1^*(\eta') & \longrightarrow & \beta_2^*(\eta').
\end{array}
\]
De cette façon, les objets de $\mathcal F_{S'}$ munis d'une donnée de
recollement (relativement à $\beta_1,\beta_2$) forment alors une catégorie. Si
$\alpha:S'\longrightarrow S$ est un morphisme tel que
$\alpha\beta_1=\alpha\beta_2$, alors pour tout $\xi\in\mathcal F_S$, l'objet
$\xi'=\alpha^*(\xi)$

\src{302}
de $\mathcal F_{S'}$ est muni d'une donnée de recollement canonique, puisque
\[
  \beta_i^*\alpha^*(\xi)\simeq(\alpha\beta_i)^*(\xi)=\gamma^*(\xi),
\]
en posant encore $\gamma=\alpha\beta_1=\alpha\beta_2$ ; de plus, si
$u:\xi\longrightarrow\eta$ est un morphisme dans $\mathcal F_S$, alors
\[
  \alpha^*(u):\alpha^*(\xi)\longrightarrow\alpha^*(\eta)
\]
est un morphisme dans $\mathcal F_{S'}$ compatible avec les données de
recollement canoniques. On obtient ainsi un foncteur canonique de la catégorie
$\mathcal F_S$ dans la catégorie des objets de $\mathcal F_{S'}$ munis d'une
donnée de recollement relativement au couple $(\beta_1,\beta_2)$. Ceci dit, on
peut aussi exprimer la définition 1.3 en disant que le diagramme (+) de ladite
définition est $\mathcal F$-exact si le foncteur précédent est « pleinement
fidèle », i.e. définit une équivalence de la catégorie $\mathcal F_S$ avec une
sous-catégorie de la catégorie des objets de $\mathcal F_{S'}$ munis d'une
donnée de recollement relativement à $(\beta_1,\beta_2)$.

\result{DÉFINITION}{1.5} On dit qu'une donnée de recollement sur
$\xi'\in\mathcal F_{S'}$ est \emph{effective} (relativement à $\alpha$) si
$\xi'$, muni de cette donnée, est isomorphe à un $\alpha^*(\xi)$, avec
$\xi\in\mathcal F_S$.

Dans le cas où le diagramme (+) est $\mathcal F$-exact, l'objet $\xi$ de la
définition précédente est donc déterminé à un isomorphisme unique près, et la
catégorie $\mathcal F_S$ est équivalente à la catégorie des objets de
$\mathcal F_{S'}$ munis d'une donnée de recollement effective.

\medskip\noindent
c. Le cas le plus important est celui où
\[
  S''=S'\times_S S',
\]
les $\beta_i$ étant les deux projections $p_1$ et $p_2$ de
$S'\times_S S'$ sur ses deux facteurs (on suppose maintenant que dans
$\mathcal C$ les produits fibrés existent). On a alors deux conditions
nécessaires immédiates pour qu'une donnée de recollement
$\varphi:p_1^*(\xi')\longrightarrow p_2^*(\xi')$ sur un
$\xi'\in\mathcal F_{S'}$ soit effective. On a
\[
  \tag{i}
  \Delta^*(\varphi)=\operatorname{id}_{\xi'},
\]
où $\Delta:S'\longrightarrow S'\times_S S'$ désigne le morphisme diagonal, et
où on a identifié $\Delta^*p_i^*(\xi')$ à
$(p_i\Delta)^*(\xi')=\xi'$. On a
\[
  \tag{ii}
  p_{31}^*(\varphi)=p_{32}^*(\varphi)\,p_{21}^*(\varphi),
\]

\src{303}
où $p_{ij}$ désigne la projection canonique de
$S'\times_S S'\times_S S'$ sur le produit partiel de son $i$-ième et
$j$-ième facteur.

\result{DÉFINITION}{1.6} On appelle \emph{donnée de descente} sur
$\xi'\in\mathcal F_{S'}$, relativement au morphisme
$\alpha:S'\longrightarrow S$, une donnée de recollement sur $\xi'$
relativement au couple $(p_1,p_2)$ des projections canoniques
$S'\times_S S'\longrightarrow S'$, satisfaisant aux conditions (i) et (ii)
ci-dessus.

\result{DÉFINITION}{1.7} Un morphisme $\alpha:S'\longrightarrow S$ est appelé
un \emph{morphisme de $\mathcal F$-descente} si le diagramme de morphismes
\[
  S\xleftarrow{\alpha}S'
  \overset{p_1}{\underset{p_2}{\rightrightarrows}}S'\times_S S'
\]
est $\mathcal F$-exact (définition 1.3) ; on dit que $\alpha$ est un
\emph{morphisme de $\mathcal F$-descente strict} si de plus toute donnée de
descente (définition 1.6) sur un objet de $\mathcal F_{S'}$ est effective.

Cette dernière condition peut aussi, de façon plus imagée, s'énoncer ainsi :
« Il revient au même » de se donner un objet de $\mathcal F_S$, ou un objet de
$\mathcal F_{S'}$ muni d'une donnée de descente.

Notons que si un morphisme $\alpha:S'\longrightarrow S$ admet une « section »
$s:S\longrightarrow S'$ (i.e. un morphisme $s$ tel que
$\alpha s=\operatorname{id}_S$), alors c'est un morphisme de
$\mathcal F$-descente strict.\footnote{L'erratum de 1962 précise qu'il est
inutile de supposer au préalable que $\alpha$ soit un morphisme de
$\mathcal F$-descente.} Si $\xi'\in\mathcal F_{S'}$ est muni d'une donnée de
descente relativement à $\alpha$, il « provient » en effet de
$\xi=s^*(\xi')$.

\medskip\noindent
d. On peut présenter de façon plus intuitive les notions précédentes, en
introduisant, pour un objet-paramètre $T$ de $\mathcal C$ au-dessus de $S$,
l'ensemble
\[
  \operatorname{Hom}_S(T,S')=S'(T),
\]
dont les éléments seront désignés par $t,t'$, etc. Un objet
$\xi'\in\mathcal F_{S'}$ étant donné, à tout $t\in S'(T)$ correspond alors un
objet $t^*(\xi')$ de $\mathcal F_T$, qui sera aussi noté $\xi'_t$. Une donnée de
recollement sur $\xi'$ (relativement à $p_1,p_2$) est alors définie par la
donnée, pour tout $T$ sur $S$, et tout couple de points
$t,t'\in S'(T)$, d'un isomorphisme
\[
  \varphi_{t',t}:\xi'_t\longrightarrow\xi'_{t'}
\]
(satisfaisant des conditions évidentes de fonctorialité, pour $T$ variable).
Les conditions (i) et (ii) de 1.c s'écrivent alors
\[
  \tag{i bis}
  \varphi_{t,t}=\operatorname{id}_{\xi'_t}
  \qquad\text{pour tout $T$ et tout $t\in S'(T)$},
\]

\src{304}
\[
  \tag{ii bis}
  \varphi_{t,t''}=\varphi_{t,t'}\varphi_{t',t''}
  \qquad\text{pour tout $T$ et tous $t,t',t''\in S'(T)$}.
\]
On voit d'ailleurs que (ii bis), en faisant $t=t'=t''$, implique
$\varphi_{t,t}=\varphi_{t,t}^2$, d'où, puisque $\varphi_{t,t}$ est un
isomorphisme par hypothèse, la relation (i bis), qui est donc en fait une
conséquence de (ii bis) (donc (i) est une conséquence de (ii)). Mais si on ne
suppose plus a priori les $\varphi_{t,t'}$ des isomorphismes (i.e. que
$\varphi:p_1^*(\xi')\longrightarrow p_2^*(\xi')$ est un isomorphisme), alors
(ii bis) n'implique plus nécessairement (i bis) ; la conjonction de (ii bis) et
(i bis) implique cependant que les $\varphi_{t,t'}$ sont des isomorphismes (car
on aura
$\varphi_{t,t'}\varphi_{t',t}=\varphi_{t,t}=\operatorname{id}_{\xi'_t}$).

\subsection*{2. Diagrammes exacts et épimorphismes stricts, morphismes de
descente. Exemples}
\addcontentsline{toc}{subsection}{2. Diagrammes exacts et épimorphismes stricts}

\result{a. DÉFINITION}{2.1} Soit $\mathcal C$ une catégorie. Un diagramme
\[
  T\xrightarrow{\alpha}T'
  \overset{\beta_1}{\underset{\beta_2}{\rightrightarrows}}T''
\]
de morphismes dans $\mathcal C$ est dit \emph{exact} si, pour tout
$Z\in\mathcal C$, le diagramme correspondant d'applications ensemblistes
\[
  \operatorname{Hom}(Z,T)\longrightarrow
  \operatorname{Hom}(Z,T')
  \rightrightarrows\operatorname{Hom}(Z,T'')
\]
est exact (définition 1.2). On dit alors que $(T,\alpha)$ (ou, par abus de
langage, que $T$) est un \emph{noyau} du couple de morphismes
$(\beta_1,\beta_2)$.

Ce noyau est évidemment déterminé à isomorphisme unique près. Si $\mathcal C$
est la catégorie des ensembles, la définition précédente est compatible avec
la définition 1.2. On définit de façon duale l'exactitude d'un diagramme
\[
  S\xleftarrow{\alpha}S'
  \overset{\beta_1}{\underset{\beta_2}{\rightrightarrows}}S''
\]
de morphismes dans $\mathcal C$ ; on dit alors que $(S,\alpha)$ est un
\emph{conoyau} du couple de morphismes $(\beta_1,\beta_2)$.

\result{DÉFINITION}{2.2} Un morphisme $\alpha:S'\longrightarrow S$ est appelé
un \emph{épimorphisme strict} si c'est un épimorphisme, et si, pour tout
morphisme $u:S'\longrightarrow Z$, la condition nécessaire suivante est aussi
suffisante pour que $u$ se factorise en $S'\longrightarrow S\longrightarrow Z$ :
pour tout $S''\in\mathcal C$ et tout couple de morphismes
$\beta_1,\beta_2:S''\longrightarrow S'$ tel que
$\alpha\beta_1=\alpha\beta_2$, on a aussi $u\beta_1=u\beta_2$.

Si le produit fibré $S'\times_S S'$ existe, il revient au même de dire que le
diagramme
\[
  S\xleftarrow{\alpha}S'
  \overset{p_1}{\underset{p_2}{\rightrightarrows}}S'\times_S S'
\]

\src{305}
est exact, i.e. que $S$ est un conoyau du couple $(p_1,p_2)$. En tout cas, un
morphisme conoyau est un épimorphisme strict. Notons aussi qu'un épimorphisme
strict qui est un monomorphisme est un isomorphisme. Nous laissons au lecteur
le soin de développer la notion duale de \emph{monomorphisme strict}.

Pour préciser les relations entre la notion de morphisme de
$\mathcal F$-descente et la notion d'épimorphisme strict, nous introduisons
encore les définitions suivantes :

\result{DÉFINITION}{2.3} Un morphisme $\alpha:S'\longrightarrow S$ est appelé
un \emph{épimorphisme universel} (resp. un \emph{épimorphisme strict universel})
si, pour tout $T$ sur $S$, le produit fibré $T'=S'\times_ST$ existe, et la
projection $T'\longrightarrow T$ est un épimorphisme (resp. un épimorphisme
strict).

Dans les très bonnes catégories (telles la catégorie des ensembles, la
catégorie des faisceaux d'ensembles sur un espace topologique, les catégories
abéliennes, etc.) les quatre notions d'épijectivité ainsi introduites
coïncident ; elles sont au contraire toutes distinctes dans une catégorie telle
que la catégorie des préschémas, ou la catégorie des préschémas au-dessus d'un
préschéma non vide donné $S$, même si on se borne à des $S$-schémas finis sur
$S$.

\result{DÉFINITION}{2.4} Un morphisme $\alpha:S'\longrightarrow S$ est appelé
\emph{morphisme de descente} (resp. un \emph{morphisme de descente strict}) si
c'est un morphisme de $\mathcal F$-descente (resp. un morphisme de
$\mathcal F$-descente strict) (cf. définition 1.7), où $\mathcal F$ désigne la
catégorie fibrée de base $\mathcal C$ des objets de $\mathcal C$ sur des objets
de $\mathcal C$ (n\textsuperscript{o} 1, exemple 1).

\result{PROPOSITION}{2.1} Si dans $\mathcal C$ les produits finis et les
produits fibrés (finis) existent, alors il y a identité dans $\mathcal C$ entre
morphismes de descente et épimorphismes stricts universels.

\result{b. EXEMPLES}{} Soit $\mathcal C$ la catégorie des préschémas. Soit
$S\in\mathcal C$, soient $S'$ et $S''$ deux préschémas finis sur $S$, i.e.
correspondant à des faisceaux d'algèbres $\mathcal A'$, $\mathcal A''$ sur $S$
qui, en tant que faisceaux de modules, sont quasi-cohérents et de type fini
(i.e. cohérents si $S$ est localement noethérien). Soit
$\alpha:S'\longrightarrow S$ le morphisme structural de $S'$, et soient
$\beta_1,\beta_2$ deux $S$-morphismes de $S''$ dans $S'$, définis par des
homomorphismes d'algèbres $\mathcal A'\longrightarrow\mathcal A''$ désignés
encore par $\beta_1,\beta_2$. Utilisant le fait qu'un morphisme fini est fermé
(premier théorème de Cohen--Seidenberg), on prouve facilement que le diagramme
\[
  \tag{+}
  S\xleftarrow{\alpha}S'
  \overset{\beta_1}{\underset{\beta_2}{\rightrightarrows}}S''
\]
dans $\mathcal C$ est exact si et seulement si le diagramme de faisceaux

\src{306}
\[
  \mathcal O_S=\mathcal A\longrightarrow\mathcal A'
  \overset{\beta_1}{\underset{\beta_2}{\rightrightarrows}}\mathcal A''
\]
sur $S$ est exact. En particulier, si $\alpha:S'\longrightarrow S$ est un
morphisme fini correspondant à un faisceau $\mathcal A'$ d'algèbres sur $S$,
alors $\alpha$ est un épimorphisme strict si et seulement si le diagramme de
faisceaux
\[
  \mathcal O_S=\mathcal A\longrightarrow\mathcal A'
  \overset{p_1}{\underset{p_2}{\rightrightarrows}}
  \mathcal A'\otimes_{\mathcal A}\mathcal A'
\]
est exact (c'est un épimorphisme si et seulement si
$\mathcal A\longrightarrow\mathcal A'$ est injectif). Si $S$ est affine
d'anneau $A$, donc $S'$ affine d'anneau $A'$ fini sur $A$, alors
$S'\longrightarrow S$ est un épimorphisme strict si et seulement si
$A\longrightarrow A'$ est un isomorphisme de $A$ sur le sous-anneau de $A'$
formé des $x'\in A'$ tels que
\[
  1_{A'}\otimes_Ax'-x'\otimes_A1_{A'}=0
\]
(c'est un épimorphisme si et seulement si $A\longrightarrow A'$ est injectif).
Comme nous l'avons déjà signalé, même si $S$ est le schéma d'un anneau local
artinien, un morphisme fini $S'\longrightarrow S$ qui est un épimorphisme n'est
pas nécessairement un épimorphisme strict. Cependant, on peut prouver que, si
$S$ est un préschéma noethérien, tout morphisme fini $S'\longrightarrow S$ qui
est un épimorphisme est le composé d'une suite finie d'épimorphismes stricts
(également finis). Cela montre d'ailleurs que le composé de deux épimorphismes
stricts n'est pas nécessairement un épimorphisme strict.

\result{c}{} Si $(+)$ est un diagramme exact de morphismes finis, alors, pour
tout morphisme plat $T\longrightarrow S$ de préschémas, le diagramme transformé
de $(+)$ par le changement de base $T\longrightarrow S$ est encore exact. Il en
résulte que, si $X,Y$ sont deux $S$-préschémas, $X$ étant plat sur $S$, alors le
diagramme d'applications ensemblistes suivant (où $X',Y'$ sont les images
réciproques de $X,Y$ sur $S'$, et $X'',Y''$ leurs images réciproques sur $S''$)
\[
  \operatorname{Hom}_S(X,Y)\longrightarrow
  \operatorname{Hom}_{S'}(X',Y')
  \rightrightarrows
  \operatorname{Hom}_{S''}(X'',Y'')
\]
est exact. En particulier, si $\mathcal F$ désigne la catégorie fibrée de base
la catégorie $\mathcal C$ des préschémas, telle que, pour $X\in\mathcal C$,
$\mathcal F_X$ soit la catégorie des $X$-préschémas plats, alors le diagramme
$(+)$ est $\mathcal F$-exact. (Ce résultat devient faux si on ne fait pas
d'hypothèse de platitude ; en particulier un épimorphisme strict fini n'est pas
nécessairement un morphisme de descente.) On voit de même que $(+)$ est
$\mathcal F$-exact si $\mathcal F$ désigne la catégorie fibrée pour laquelle
$\mathcal F_X$ est la catégorie des faisceaux quasi-cohérents et plats sur le
préschéma $X$ (ici encore l'hypothèse

\src{307}
de platitude est essentielle). Dans l'un et l'autre cas, la question de
l'effectivité d'une donnée de recollement (et plus particulièrement d'une
donnée de descente, lorsque $S''=S'\times_SS'$) sur un objet plat au-dessus de
$S'$ est délicate (et sa réponse dans divers cas particuliers est un des objets
principaux des présents exposés). Le conférencier ignore si, pour tout
épimorphisme strict fini $S'\longrightarrow S$, toute donnée de descente sur un
faisceau quasi-cohérent plat sur $S'$ est effective (même en supposant que $S$
est le spectre d'un anneau local artinien, et en se bornant aux faisceaux
localement libres de rang 1). Plus généralement, soient $A$ un anneau, $A'$ une
$A$-algèbre (tout est commutatif) telle que le diagramme d'applications
\[
  A\longrightarrow A'\rightrightarrows A'\otimes_AA'
\]
soit exact, ce qui équivaut aussi au fait que le morphisme
$S'\longrightarrow S$ correspondant pour les spectres de $A,A'$ soit un
morphisme de $\mathcal F$-descente, où $\mathcal F$ est la catégorie fibrée des
faisceaux quasi-cohérents plats. Soit $M'$ un $A'$-module plat muni d'une donnée
de descente à $A$, i.e. d'un isomorphisme
\[
  \varphi:M'\otimes_AA'\xrightarrow{\sim}A'\otimes_AM'
\]
de $A'\otimes_AA'$-modules, satisfaisant les conditions (i) et (ii) de 1.c (que
nous laissons au lecteur le soin d'expliciter en termes de modules). Cette
donnée est-elle effective (relativement à la catégorie fibrée des faisceaux
quasi-cohérents plats) ? Soit $M$ le sous-ensemble de $M'$ formé des $x'\in M'$
tels que
\[
  \varphi(x'\otimes_A1_{A'})=1_{A'}\otimes_Ax',
\]
c'est un sous-$A$-module de $M'$. L'injection canonique
$M\longrightarrow M'$ définit un homomorphisme de $A'$-modules
$M\otimes_AA'\longrightarrow M'$. L'effectivité de $\varphi$ signifie alors
ceci : $M$ est un $A$-module plat, et l'homomorphisme précédent est un
isomorphisme.

\result{REMARQUE}{} Dans les considérations précédentes, nous n'avions fait
aucune hypothèse de platitude sur les morphismes du diagramme $(+)$, ce qui nous
obligeait, pour avoir une technique de descente, à faire des hypothèses de
platitude sur les objets au-dessus de $S,S'$ qu'on considère. Dans le paragraphe
2, nous ferons une hypothèse de platitude sur $\alpha:S'\longrightarrow S$, ce
qui nous permettra d'avoir une technique de descente pour des objets au-dessus
de $S,S'$ qui ne sont plus soumis à aucune condition de platitude. Dans tous
les cas, il y a une hypothèse de platitude qui intervient. C'est là une des
principales raisons de l'importance de la notion de platitude en géométrie
algébrique (dont le rôle ne pouvait apparaître tant qu'on se bornait à des
corps de base, sur lesquels n'importe quoi, en effet, est plat !).

\src{308}
\subsection*{3. Application aux étalements}
\addcontentsline{toc}{subsection}{3. Application aux étalements}

Soient $A$ un anneau local, $B$ une algèbre locale sur $A$ dont l'idéal maximal
induit celui de $A$. Nous dirons que $B$ est \emph{étalé} sur $A$ (au lieu de
« non ramifié », utilisé par ailleurs) s'il satisfait les conditions suivantes :
\begin{enumerate}
\item[(i)] $B$ est plat sur $A$ ;
\item[(ii)] $B/\mathfrak mB$ est une extension finie séparable de
$A/\mathfrak m=k$ (où $\mathfrak m$ désigne l'idéal maximal de $A$).
\end{enumerate}
Lorsque $A$ et $B$ sont noethériens et $k$ algébriquement clos, cela signifie
que l'homomorphisme $\widehat A\longrightarrow\widehat B$ sur les complétés qui
prolonge $A\longrightarrow B$ est un isomorphisme. Un morphisme de type fini
$f:T\longrightarrow S$ est dit \emph{étale en $x\in T$}, ou encore $T$ est dit
\emph{étalé sur $S$ en $x$}, si $\mathcal O_x$ est étalé sur
$\mathcal O_{f(x)}$, et $f$ est dit \emph{étale}, ou encore $f$ est appelé un
\emph{étalement}, ou $T$ est dit \emph{étalé sur $S$}, si $f$ est étale en tout
$x\in T$. Remarquons d'ailleurs que, si $S$ est localement noethérien,
l'ensemble des points de $T$ où $f$ est étale est ouvert, et l'utilisation du
« main theorem » de Zariski permet de préciser la structure de $T/S$ au
voisinage d'un tel point (par une équation de type bien connu).

Lorsque $S$ est un schéma de type fini sur le corps des complexes, il lui
correspond un espace analytique $\overline S$ au sens de Serre [5], à cela près
que $\overline S$ peut avoir des éléments nilpotents dans son faisceau
structural, ce qui ne change rien d'essentiel dans [5]. On voit alors facilement
que $f$ est un étalement si et seulement si
$\overline f:\overline T\longrightarrow\overline S$ l'est, i.e. si tout point
de $\overline T$ admet un voisinage sur lequel $\overline f$ induit un
isomorphisme sur un ouvert de $\overline S$. En particulier, à tout
\emph{revêtement étale} $T$ de $S$ (i.e. un morphisme fini étale
$f:T\longrightarrow S$) correspond un revêtement étale $\overline T$ de
$\overline S$, qui est connexe si et seulement si $T$ l'est [5]. On voit
d'ailleurs facilement que, si $T,T'$ sont deux schémas étalés sur $S$, alors
l'application naturelle
\[
  \operatorname{Hom}_S(T,T')\longrightarrow
  \operatorname{Hom}_{\overline S}(\overline T,\overline{T'})
\]
est bijective, i.e. le foncteur $T\longmapsto\overline T$ de la catégorie des
schémas étalés sur $S$ dans la catégorie des espaces analytiques étalés sur
$\overline S$ est « pleinement fidèle », donc définit une équivalence de la
première catégorie avec une sous-catégorie de la seconde. Un théorème de
Grauert--Remmert [2] implique que, si $S$ est normal, on obtient ainsi une
équivalence de la catégorie des \emph{revêtements étales} de $S$ et de la
catégorie des revêtements étales (finis) de $\overline S$, i.e. que tout
revêtement étale de $\overline S$ est $\overline S$-isomorphe à un
$\overline T$, où $T$ est un revêtement étale de $S$. Montrons que le
\emph{théorème de Grauert--Remmert reste valable sans hypothèse de normalité sur
$S$}. Soit en effet d'abord $S'\longrightarrow S$ un épimorphisme strict fini,
supposons le

\src{309}
théorème démontré pour $S'$, montrons qu'il sera vrai pour $S$. En effet, soit
$\mathfrak T$ un revêtement étale de $\overline S$, considérons son image
réciproque $\mathfrak T'$ sur $\overline{S'}$, qui correspond à un faisceau
analytique cohérent $\mathcal A'$ d'algèbres sur $\overline{S'}$, image
réciproque du faisceau d'algèbres $\mathcal A$ sur $\overline S$ définissant
$\mathfrak T$. Par hypothèse, sur $\overline{S'}$, $\mathfrak T'$ provient d'un
revêtement étale $T'$ de $S'$, i.e. $\mathcal A'$ provient d'un faisceau
cohérent d'algèbres $A'$ sur $S'$. D'autre part, $\mathcal A'$ est muni d'une
donnée de descente canonique relativement à
$\overline{S'}\longrightarrow\overline S$, i.e. d'un isomorphisme entre ses
deux images réciproques sur
\[
  \overline{S'}\times_{\overline S}\overline{S'}
  =\overline{(S'\times_SS')},
\]
(satisfaisant aux conditions (i), (ii)), et cet isomorphisme provient, d'après
ce qui a été dit, d'un isomorphisme sur les faisceaux algébriques
correspondants, i.e. d'une donnée de descente sur $A'$ relativement à
$S'\longrightarrow S$. On vérifie facilement que cette dernière est effective
(car celle sur $\mathcal A'$ l'est, et l'effectivité d'une donnée de descente,
telle qu'elle a été explicitée au numéro précédent, se reconnaît localement sur
les \emph{complétés} des modules qui entrent en jeu). D'où un faisceau cohérent
d'algèbres $A$ sur $S$, définissant un revêtement $T$ de $S$, qui est le
revêtement cherché. Le résultat précédent reste alors manifestement valable si
$S'\longrightarrow S$ est seulement un composé d'un nombre fini
d'épimorphismes stricts finis, i.e. est un épimorphisme fini quelconque (d'après
le résultat de factorisation signalé au paragraphe 2). Il s'ensuit que le
théorème de Grauert--Remmert reste valable si $S$ est un schéma \emph{réduit},
i.e. tel que $\mathcal O_S$ n'ait pas d'éléments nilpotents, comme on voit en
introduisant son normalisé $S'$. On passe facilement de là au cas général.

Une démonstration toute analogue, utilisant encore le résultat de factorisation
pour les épimorphismes stricts finis et la nature « formelle » de l'effectivité
de données de descente, permet de prouver le résultat suivant : soit $S$ un
préschéma localement noethérien et soit $S'\longrightarrow S$ un morphisme
\emph{fini, surjectif, radiciel} (ou, ce qui revient au même, un morphisme de
type fini tel que, pour tout $T$ sur $S$, le morphisme
$T'=S'\times_ST\longrightarrow T$ soit un homéomorphisme, ce qu'on exprime
encore en disant que $S'\longrightarrow S$ est un
\emph{homéomorphisme universel}). Pour tout $T$ étalé sur $S$, considérons son
image inverse $T'=T\times_SS'$, qui est étalé sur $S'$. Alors le foncteur
$T\longmapsto T'$ est une \emph{équivalence de catégories de la catégorie des
préschémas $T$ étalés sur $S$ avec la catégorie des préschémas $T'$ étalés sur
$S'$}. (On utilise la bijectivité de
\[
  \operatorname{Hom}_S(T_1,T_2)\longrightarrow
  \operatorname{Hom}_{S'}(T'_1,T'_2)
\]
pour deux préschémas $T_1,T_2$ étalés sur $S$, fait dont la vérification
directe est facile, et le fait que le théorème énoncé est vrai si
$S'=(S,\mathcal O_S/\mathcal J)$,

\src{310}
où $\mathcal J$ est un faisceau cohérent nilpotent d'idéaux de $\mathcal O_S$
([4], lemme 6)). Noter d'ailleurs que nous ne supposons pas ici les $T$
envisagés finis sur $S$. Ce résultat implique en particulier que le morphisme
$S'\longrightarrow S$ induit un isomorphisme du groupe fondamental $[\,]$ de
$S'$ sur celui de $S$ (« \emph{invariance topologique du groupe fondamental
d'un préschéma} »).

\subsection*{4. Relations avec la 1-cohomologie}
\addcontentsline{toc}{subsection}{4. Relations avec la 1-cohomologie}

\medskip\noindent
a. Soit $\mathcal C$ une catégorie où le produit de deux objets existe toujours,
soit $T\in\mathcal C$. Pour tout ensemble fini $I\ne\varnothing$, on peut
considérer $T^I$ ; pour $I$ variable, on obtient ainsi un foncteur covariant de
la catégorie des ensembles finis non vides dans $\mathcal C$, i.e. ce qu'on
peut appeler un \emph{objet simplicial} de $\mathcal C$, noté $K_T$. Ce dernier
dépend de façon covariante de $T$ ; d'ailleurs, si $u,v$ sont deux morphismes
$T\longrightarrow T'$, alors les morphismes correspondants
$K_T\longrightarrow K_{T'}$ sont homotopes. Disons que $T$ \emph{domine} $T'$
si $\operatorname{Hom}(T,T')\ne\varnothing$, c'est là une relation de préordre
filtrante croissante dans $\mathcal C$. Il résulte de ce qui précède que, si $T$
domine $T'$, il existe une classe (à une homotopie près) canonique
d'homomorphismes d'objets simpliciaux $K_T\longrightarrow K_{T'}$ ; en
particulier, si $K_T$ et $K_{T'}$ sont tels que chacun domine l'autre, alors
$K_T$ et $K_{T'}$ sont homotopiquement équivalents. Soit maintenant $F$ un
foncteur (contravariant pour fixer les idées) de $\mathcal C$ dans une catégorie
abélienne $\mathcal C'$, alors
\[
  C^*(T,F)=F(K_T)
\]
est un objet cosimplicial de $\mathcal C'$, donc définit de façon bien connue
un complexe (de cochaînes) dans $\mathcal C'$, dont on peut prendre la
cohomologie
\[
  H^*(T,F)=H^*(C^*(T,F))=H^*(F(K_T))
\]
(on pourra mettre un $\mathcal C$ en indice du $H^*$ s'il y a possibilité de
confusion). C'est là un foncteur cohomologique en $F$, dont la variance pour $T$
variable résulte de ce qui a été dit sur les $K_T$ ; de façon précise, pour $F$
fixé et $T$ variable dans $\mathcal C$ (préordonné par la relation de
domination), les $H^*(T,F)$ forment un système inductif d'objets gradués de
$\mathcal C'$ ; en particulier, si $T$ et $T'$ sont tels que chacun domine
l'autre, alors $H^*(T,F)$ et $H^*(T',F)$ sont canoniquement isomorphes.

Supposons que dans $\mathcal C$ les produits fibrés existent, alors on peut,
pour $S\in\mathcal C$ fixé, appliquer ce qui précède à la catégorie
$\mathcal C_S$ des objets de $\mathcal C$ au-dessus de $S$ ; on écrira
$C^*(T/S,F)$ et $H^*(T/S,F)$ au lieu de $C^*(T,F)$ et $H^*(T,F)$ si on veut
préciser que l'on se place dans la catégorie $\mathcal C_S$ ; ainsi,

\src{311}
$C^*(T/S,F)$ est un complexe de cochaînes dans $\mathcal C'$ qui, en dimension
$n$, est égal à
\[
  F(T\times_ST\times_S\cdots\times_ST)
\]
(où la parenthèse comporte $n+1$ facteurs).

Notons que, comme d'habitude, on peut définir $H^0(T/S,F)$ sans supposer la
catégorie $\mathcal C'$ abélienne : c'est le noyau (définition 2.1), s'il
existe, du couple de morphismes $F(p_i)$ ($i=1,2$)
\[
  F(T)\rightrightarrows F(T\times_ST)
\]
correspondant aux deux projections
$p_1,p_2:T\times_ST\rightrightarrows T$. En particulier, on aura un morphisme
naturel (dit d'\emph{augmentation})
\[
  F(S)\longrightarrow H^0(T/S,F)
\]
qui sera un isomorphisme dans les cas favorables (en particulier si
$T\longrightarrow S$ est un épimorphisme strict et si $F$ est « exact à
gauche »). De même, lorsque $F$ prend ses valeurs dans la catégorie des groupes
dans une catégorie $\mathcal C''$, on peut aussi définir $H^1(T/S,F)$ ; dans le
cas où $\mathcal C''$ est la catégorie des ensembles (i.e. $F$ prend ses
valeurs dans la catégorie des groupes ordinaires, non nécessairement
commutatifs), $H^1(T/S,F)$ est le quotient du sous-groupe $Z^1(T/S,F)$ de
$C^1(T/S,F)=F(T\times_ST)$ formé des $g$ tels que
\[
  F(p_{31})(g)=F(p_{32})(g)\,F(p_{21})(g)
\]
par le groupe d'opérateurs $F(T)$, opérant sur $C^1(T/S,F)$ et, en particulier,
sur le sous-ensemble $Z^1(T/S,F)$ par
\[
  \rho(g')\cdot g=F(p_2)(g')\,g\,F(p_1)(g')^{-1}.
\]

\medskip\noindent
b. Soit par exemple $\mathcal F$ une catégorie fibrée de base $\mathcal C$.
Soient $\xi,\eta\in\mathcal F_S$, et pour tout $S'$ sur $S$, soit
\[
  F_{\xi,\eta}(S')=
  \operatorname{Hom}(\xi\times_SS',\eta\times_SS').
\]
Ainsi, $F_{\xi,\eta}$ est un foncteur contravariant de $\mathcal C_S$ dans la
catégorie des ensembles. Ceci posé, \emph{dire que le morphisme d'augmentation}
\[
  F_{\xi,\eta}(S)\longrightarrow H^0(S'/S,F_{\xi,\eta})
\]
\emph{est un isomorphisme pour tout couple d'éléments
$\xi,\eta\in\mathcal F_S$ signifie que
$\alpha:S'\longrightarrow S$ est un morphisme de $\mathcal F$-descente}
(définition 1.7).

\medskip\noindent
c. Posons de même, pour $\xi\in\mathcal F_S$ et tout objet $S'$ de $\mathcal C$
au-dessus de $S$,
\[
  G_\xi(S')=\operatorname{Aut}(\xi\times_SS') ;
\]
on a ainsi défini un foncteur contravariant $G_\xi$ de $\mathcal C_S$ dans la
catégorie des

\src{312}
groupes. Ceci posé, on constate que $Z^1(S'/S,G_\xi)$ \emph{s'identifie
canoniquement à l'ensemble des données de descente sur
$\xi'=\xi\times_SS'$ relativement à $S'\longrightarrow S$} (définition 1.6),
et $H^1(S'/S,G_\xi)$ \emph{s'identifie à l'ensemble des classes (à un
isomorphisme près) d'objets de $\mathcal F_{S'}$ munis d'une donnée de descente
relativement à $\alpha:S'\longrightarrow S$, qui, en tant qu'objets de
$\mathcal F_{S'}$, sont isomorphes à $\xi'=\xi\times_SS'$}. Si donc
$\alpha:S'\longrightarrow S$ est un morphisme de $\mathcal F$-descente
(cf. (b)), alors $H^1(S'/S,G_\xi)$ \emph{contient comme sous-ensemble
l'ensemble des classes (à un isomorphisme près) d'objets $\eta$ de
$\mathcal F_S$ tels que $\eta\times_SS'$ soit isomorphe dans
$\mathcal F_{S'}$ à $\xi\times_SS'$ ; et cette inclusion est une identité si
et seulement si toute donnée de descente sur $\xi'=\xi\times_SS'$
relativement à $\alpha:S'\longrightarrow S$ est effective}. (Ce sera le cas en
particulier si $\alpha:S'\longrightarrow S$ est un morphisme de
$\mathcal F$-descente strict.)

\result{REMARQUE}{} Les complexes de cochaînes du type $C^*(T/S,F)$ contiennent
comme cas particuliers la plupart des complexes standard connus (cohomologie
de Čech, cohomologie des groupes, etc.), et jouent un rôle important en
géométrie algébrique, notamment dans la « cohomologie de Weil » des préschémas.

\medskip\noindent
d. \result{EXEMPLE}{1} Soit $S'$ un objet au-dessus de $S\in\mathcal C$, et
soit $\Gamma$ un groupe d'automorphismes de $S'$ tel que $S'$ soit
« formellement principal sur $S$, de groupe $\Gamma$ », i.e. tel que le
morphisme naturel
\[
  \Gamma\times S'\longrightarrow S'\times_SS',
\]
où $\Gamma\times S'$ désigne la somme directe de $\Gamma$ copies de $S'$, soit
un isomorphisme. (On suppose que dans $\mathcal C$ les sommes directes qui
interviennent ici existent.) Soit $F$ un foncteur contravariant de $\mathcal C$
dans la catégorie des groupes abéliens. Alors $C^*(S'/S,F)$ est canoniquement
isomorphe au groupe simplicial de cochaînes standard homogènes
$C^*(\Gamma,F(S'))$, donc $H^*(S'/S,F)$ est canoniquement isomorphe à
$H^*(\Gamma,F(S'))$.

\medskip\noindent
e. \result{EXEMPLE}{2} Soit $\mathcal C$ la catégorie des préschémas. On désigne
par $G_a$ (« groupe additif ») le foncteur contravariant de $\mathcal C$ dans
la catégorie des groupes abéliens, défini par
\[
  G_a(X)=H^0(X,\mathcal O_X).
\]
On définit de même le foncteur $G_m$ (« groupe multiplicatif ») par
\[
  G_m(X)=H^0(X,\mathcal O_X)^*
\]
(groupe des éléments inversibles de l'anneau $H^0(X,\mathcal O_X)$), et plus
généralement le foncteur $\operatorname{Gl}(n)$ (« groupe linéaire d'ordre
$n$ ») par

\src{313}
\[
  \operatorname{Gl}(n)(X)=
  \operatorname{Gl}\bigl(n,H^0(X,\mathcal O_X)\bigr),
\]
qui est un foncteur de $\mathcal C$ dans la catégorie des groupes (non
nécessairement commutatifs si $n>1$ ; pour $n=1$ on retrouve $G_m$). On peut
d'ailleurs interpréter $\operatorname{Gl}(n)$ comme un foncteur-automorphisme
(cf. (c)) en considérant la catégorie fibrée $\mathcal F$ de base $\mathcal C$
telle que, pour $X\in\mathcal C$, $\mathcal F_X$ soit la catégorie des
faisceaux localement libres sur $X$ : on a en effet
\[
  \operatorname{Gl}(n)(X)=
  \operatorname{Aut}_{\mathcal F_X}(\mathcal O_X^n).
\]
D'après (b), il s'ensuit que, si $\alpha:S'\longrightarrow S$ est un morphisme
de $\mathcal F$-descente (cf. paragraphe 2.c), $H^1(S'/S,\operatorname{Gl}(n))$
contient l'ensemble des classes (à un isomorphisme près) de faisceaux localement
libres sur $S$ dont l'image inverse sur $S'$ est isomorphe à
$\mathcal O_{S'}^n$, et cette inclusion est une égalité si et seulement si
toute donnée de descente sur $\mathcal O_{S'}^n$ (relativement à
$\alpha:S'\longrightarrow S$) est effective. Lorsque $S$ est le spectre d'un
anneau local, cela signifie donc
$H^1(S'/S,\operatorname{Gl}(n))=(e)$, puisque tout faisceau localement libre
sur $S$ est alors trivial.

Notons l'équivalence des conditions suivantes sur un morphisme
$\alpha:S'\longrightarrow S$ :
\begin{enumerate}
\item[(i)] l'homomorphisme d'augmentation
\[
  H^0(S,\mathcal O_S)=G_a(S)\longrightarrow H^0(S'/S,G_a)
\]
est un isomorphisme ;
\item[(ii)] $S'\longrightarrow S$ est un morphisme de $\mathcal F$-descente
($\mathcal F$ étant la catégorie fibrée de base $\mathcal C$ envisagée
ci-dessus).
\end{enumerate}
Si $S'\longrightarrow S$ est fini, ces conditions équivalent aussi à :
\begin{enumerate}
\item[(iii)] $S'\longrightarrow S$ est un épimorphisme strict
(cf. paragraphe 2.c).
\end{enumerate}

Supposons maintenant que $S=\operatorname{Spec}(A)$,
$S'=\operatorname{Spec}(A')$ ; alors on a
\[
  C^n(S'/S,G_a)=C^n(A'/A,G_a)=\bigotimes_A^{n+1}A',
\]
l'opérateur cobord
$C^n(A'/A,G_a)\longrightarrow C^{n+1}(A'/A,G_a)$ étant somme alternée des
opérateurs faces
\[
  \partial_i(x_0\otimes x_1\otimes\cdots\otimes x_n)
  =x_0\otimes\cdots\otimes x_{i-1}\otimes1_{A'}\otimes
   x_i\otimes\cdots\otimes x_n.
\]
De même,
\[
  C^n(S'/S,G_m)=C^n(A'/A,G_m)
  \simeq\left(\bigotimes_A^{n+1}A'\right)^*,
\]
les opérations simpliciales dans $C^*(A'/A,G_m)$ étant induites par celles de
$C^*(S'/S,G_a)$. On explicite de même les opérations simpliciales dans
$C^*(A'/A,\operatorname{Gl}(n))$. \emph{Dans tous les cas à la connaissance du
conférencier, on a}
\[
  H^i(A'/A,G_a)=0\qquad\text{pour $i>0$},
\]
\emph{et si $A$ est local, on a}
\[
  H^1(A'/A,G_m)=0
  \quad\text{et plus généralement}\quad
  H^1(A'/A,\operatorname{Gl}(n))=(e)
\]
(lorsque $S'\longrightarrow S$ est un morphisme de $\mathcal F$-descente, i.e.
le diagramme
$A\longrightarrow A'\rightrightarrows A'\otimes_AA'$ est exact, comparer avec
le paragraphe 2.c). On notera que \emph{le « théorème 90 » de Hilbert n'est
autre que la relation}
\src{314}
\[
  H^1(S'/S,G_m)=0
\]
\emph{lorsque $A$ est un corps et $A'$ une extension galoisienne finie de ce
dernier (cf. exemple 1), et peut encore s'exprimer en disant que, dans le cas
envisagé, $S'\longrightarrow S$ est un morphisme de descente strict pour la
catégorie fibrée des faisceaux localement libres de rang 1.} C'est sous cette
dernière forme qu'il convient de généraliser le théorème de Hilbert, en variant
les hypothèses aussi bien sur le morphisme $S'\longrightarrow S$ que sur les
faisceaux quasi-cohérents envisagés.

Notons enfin l'équivalence des propriétés suivantes, lorsque $A$ est un anneau
local artinien d'idéal maximal $\mathfrak m$, $A'$ une $A$-algèbre (en
désignant, pour tout entier $k>0$, par $A_k$ (resp. $A'_k$) les anneaux
$A/\mathfrak m^{k+1}$ (resp. $A'/\mathfrak m^{k+1}A'$)) :
\begin{enumerate}
\item[(i)] $H^1(A'_k/A_k,G_a)=0$ pour tout $k$.
\item[(ii)] $H^1(A'_k/A_k,G_m)=0$ pour tout $k$.
\item[(iii)] $H^1(A'_k/A_k,\operatorname{Gl}(n))=(e)$ pour tout $k$ et tout
$n$.
\end{enumerate}
Si $S'\longrightarrow S$ est un épimorphisme strict, alors les conditions
précédentes impliquent même que c'est un morphisme de descente strict pour les
modules libres (de type fini ou non) sur $A'$.

\result{REMARQUE}{} La définition des groupes $H^i(S'/S,G_m)$, dans le cas où
$S,S'$ sont des schémas de corps $A,A'$, est due à Amitsur. Le groupe
$H^2(S'/S,G_m)$ est particulièrement intéressant comme variante « globale »
du groupe de Brauer, variante pour laquelle on pourra se référer à [1],
chapitre VII.

\fsection{B}{Descente par morphismes fidèlement plats}

\subsection*{1. Énoncé des théorèmes de descente}
\addcontentsline{toc}{subsection}{1. Énoncé des théorèmes de descente}

\result{DÉFINITION}{1.1} Un morphisme
$\alpha:S'\longrightarrow S$ de préschémas est dit \emph{plat} si, pour tout
$x'\in S'$, $\mathcal O_{x'}$ est un module plat sur l'anneau
$\mathcal O_{\alpha(x')}$ (i.e.
$\mathcal O_{x'}\otimes_{\mathcal O_{\alpha(x')}}M$ est un foncteur exact en
le $\mathcal O_{\alpha(x')}$-module $M$). Un morphisme est dit
\emph{fidèlement plat} s'il est plat et surjectif.

Par exemple, si $S=\operatorname{Spec}(A)$,
$S'=\operatorname{Spec}(A')$, alors $S'$ est plat sur $S$ si et seulement si
$A'$ est un $A$-module plat, et $S'$ est fidèlement plat sur $S$ si et
seulement si $A'$ est un $A$-module fidèlement plat (i.e. le foncteur
$A'\otimes_A M$ en le $A$-module $M$ est exact et fidèle) ; cela signifie
aussi, dans la terminologie de Serre [5], que le couple $(A,A')$ est plat. Si
$S'$ est fidèlement plat sur $S$, alors le foncteur image inverse de faisceaux
quasi-cohérents sur $S$ est exact et fidèle, en d'autres termes, pour qu'une
suite d'homomorphismes de faisceaux

\src{315}
quasi-cohérents sur $S$ soit exacte, il faut et il suffit que son image
réciproque sur $S'$ le soit (en particulier, pour qu'un homomorphisme de
faisceaux quasi-cohérents sur $S$ soit un monomorphisme, resp. un
épimorphisme, resp. un isomorphisme, il faut et il suffit que son image inverse
sur $S'$ le soit). Cette propriété reste vraie si on se restreint au-dessus
d'un ouvert quelconque de $S$, et sous cette forme caractérise les morphismes
fidèlement plats.

\result{DÉFINITION}{1.2} Un morphisme $\alpha:S'\longrightarrow S$ est dit
\emph{quasi-compact} si l'image inverse de toute partie ouverte quasi-compacte
$U$ de $S$ est quasi-compacte (i.e. réunion finie d'ouverts affines).

Il suffit évidemment de vérifier cette propriété pour les ouverts affines de
$S$. Par exemple, un morphisme affine (i.e. tel que l'image inverse d'un ouvert
affine soit affine) est quasi-compact.

La classe des morphismes plats, resp. fidèlement plats, resp. quasi-compacts,
est stable par composition et par « extension de la base », et contient bien
entendu les isomorphismes.

\result{THÉORÈME}{1} Soit $\alpha:S'\longrightarrow S$ un morphisme de
préschémas, fidèlement plat et quasi-compact. Alors $\alpha$ est un
\emph{morphisme de descente strict} (A, définition 1.7) pour la catégorie
fibrée $\mathcal F$ des faisceaux quasi-cohérents (A, paragraphe 1, exemple 2).

Cet énoncé signifie deux choses :
\begin{enumerate}
\item[(i)] Si $F$ et $G$ sont deux faisceaux quasi-cohérents sur $S$, $F'$ et
$G'$ leurs images inverses sur $S'$, alors l'homomorphisme naturel
\[
  \operatorname{Hom}(F,G)\longrightarrow\operatorname{Hom}(F',G')
\]
est une bijection du premier membre sur le sous-groupe du second formé des
homomorphismes $F'\longrightarrow G'$ qui sont compatibles avec les données de
descente canoniques sur ces faisceaux, i.e. dont les images inverses par les
deux projections de $S''=S'\times_SS'$ sur $S'$ donnent un même homomorphisme
$F''\longrightarrow G''$.
\item[(ii)] Tout faisceau quasi-cohérent $F'$ sur $S'$, muni d'une donnée de
descente relativement au morphisme $\alpha:S'\longrightarrow S$
(A, définition 1.6), est isomorphe (muni de cette donnée) à l'image inverse
d'un faisceau quasi-cohérent $F$ sur $S$.
\end{enumerate}
Faisant $F=\mathcal O_S$ dans (i), on trouve :

\result{COROLLAIRE}{1} Soit $G$ un faisceau quasi-cohérent sur $S$, soient
$G'$ et $G''$ ses images inverses sur $S'$ et sur
$S''=S'\times_SS'$, soient $p_1,p_2$ les deux projections de $S''$ sur $S'$,
alors le diagramme suivant d'applications d'ensembles

\src{316}
\[
  \Gamma(G)\xrightarrow{\alpha^*}\Gamma(G')
  \underset{p_2^*}{\overset{p_1^*}{\rightrightarrows}}\Gamma(G'')
\]
est exact (A, définition 1.(a)).

D'autre part, la conjonction de (i), (ii) et de la définition 1.1 donne le

\result{COROLLAIRE}{2} Soit $G$ comme dans le corollaire 1. Alors il y a
correspondance biunivoque entre les sous-faisceaux quasi-cohérents de $G$ et
les sous-faisceaux quasi-cohérents de $G'$ dont les images inverses sur $S''$,
par les deux projections $p_1,p_2$, donnent le même sous-faisceau de $G''$.

Bien entendu, on a un énoncé équivalent en termes de faisceaux quotients. Comme
on l'a vu (A, paragraphe 4, (e)), le théorème 1 doit être considéré comme une
généralisation du « théorème 90 » de Hilbert, et implique comme cas particuliers
diverses formulations en termes de 1-cohomologie. Pour la démonstration, on est
ramené facilement au cas où $S=\operatorname{Spec}(A)$,
$S'=\operatorname{Spec}(A')$, et pour (i) on se ramène facilement à prouver le
corollaire 1, i.e. l'exactitude du diagramme
\[
  M=A\otimes_A M\longrightarrow A'\otimes_A M
  \rightrightarrows A'\otimes_AA'\otimes_A M
\]
pour tout $A$-module $M$, ce qui résulte du lemme plus général :

\result{LEMME}{1.1} Soit $A'$ une $A$-algèbre fidèlement plate. Alors, pour
tout $A$-module $M$, le complexe $M$-augmenté
$C^*(A'/A,G_a)\otimes_A M$ (cf. A, paragraphe 4, (e)) est une
\emph{résolution} de $M$.

Il suffit de prouver que le complexe augmenté déduit du précédent par
extension de la base de $A$ à $A'$ satisfait les mêmes conclusions. Cela amène
à vérifier l'énoncé quand on remplace $A$ par $A'$ et $A'$ par
$A'\otimes_AA'$, donc nous ramène à un cas où il existe un homomorphisme de
$A$-algèbres $A'\longrightarrow A$ (ou, en termes géométriques, au cas où $S'$
sur $S$ admet une section). Dans ce cas, cela résulte des généralités du A,
paragraphe 4, (a). Signalons en passant le corollaire suivant, qui généralise
un énoncé bien connu en cohomologie de Galois (comparer A, paragraphe 4, (e)) :

\result{COROLLAIRE}{} Si $A'$ est fidèlement plat sur $A$, on a
$H^0(A'/A,G_a)=A$ et
\[
  H^i(A'/A,G_a)=0\qquad\text{pour $i\geq1$}.
\]

Pour prouver la partie (ii) du théorème 1, on procède comme pour (i) en se
ramenant au cas où $S'$ sur $S$ admet une section, où cela résulte de (i)
(cf. A, paragraphe 1, (c)).

On peut évidemment varier ad libitum le théorème 1 et ses corollaires en
introduisant des structures supplémentaires diverses sur les faisceaux (ou
systèmes de faisceaux) quasi-cohérents envisagés. Par exemple, la donnée sur
$S$ d'un faisceau
\src{317}
quasi-cohérent d'algèbres commutatives « équivaut » à la donnée sur $S'$ d'un
tel faisceau, muni d'une donnée de descente relativement à
$\alpha:S'\longrightarrow S$. Compte tenu de la correspondance fonctorielle
entre de tels faisceaux quasi-cohérents sur $S$ et des préschémas affines
au-dessus de $S$, on obtient la deuxième assertion du théorème suivant :

\result{THÉORÈME}{2} Soit $\alpha:S'\longrightarrow S$ comme dans le
théorème 1. Alors $\alpha$ est un \emph{morphisme de descente} (non strict en
général) (A, définition 2.4), et c'est un \emph{morphisme de descente strict}
pour la catégorie fibrée des schémas affines au-dessus de préschémas
(A, définition 1.7).

La première assertion du théorème signifie ceci : soient $X,Y$ deux
préschémas au-dessus de $S$, $X',Y'$ leurs images inverses sur $S'$ et
$X'',Y''$ leurs images inverses sur $S''=S'\times_SS'$, alors le diagramme
suivant d'applications naturelles
\[
  \operatorname{Hom}_S(X,Y)\xrightarrow{\alpha^*}
  \operatorname{Hom}_{S'}(X',Y')
  \underset{p_2^*}{\overset{p_1^*}{\rightrightarrows}}
  \operatorname{Hom}_{S''}(X'',Y'')
\]
est exact, i.e. $\alpha^*$ est une bijection de
$\operatorname{Hom}_S(X,Y)$ sur la partie de
$\operatorname{Hom}_{S'}(X',Y')$ formée des homomorphismes qui sont compatibles
avec les données de descente canoniques sur $X',Y'$ (i.e. dont les images
inverses par les deux projections de $S''$ sur $S'$ sont égales). Cela résulte
facilement du théorème 1, corollaire 1, lorsqu'on se borne à $Y$ affine sur
$S$ ; dans le cas général, il faut conjuguer le théorème 1 avec le résultat
suivant :

\result{LEMME}{1.2} Soit $\alpha:S'\longrightarrow S$ un morphisme fidèlement
plat et quasi-compact. Alors $S$ s'identifie à un \emph{espace topologique
quotient} de $S'$, i.e. toute partie $U$ de $S$ telle que
$\alpha^{-1}(U)$ soit ouverte est ouverte.

Pour compléter le théorème 2, il faut donner des critères d'effectivité pour
une donnée de descente sur un $S'$-préschéma $X'$ (dans le cas où $X'$ n'est
pas supposé affine sur $S'$). Remarquons d'abord qu'\emph{une telle donnée de
descente n'est pas nécessairement effective}, même si $S$ est le spectre d'un
corps $k$, $S'$ le spectre d'une extension quadratique $k'$ de ce dernier, et
$X'$ un schéma algébrique propre de dimension 2 sur $S'$ (comme on peut voir,
d'après Serre, en utilisant la surface non projective de Nagata). \emph{Pour
qu'une donnée de descente sur $X'/S'$ relativement à
$\alpha:S'\longrightarrow S$ (fidèlement plat et quasi-compact) soit effective,
il faut et il suffit que $X'$ soit réunion d'ouverts $X'_i$, affines sur $S'$,
qui soient « stables » par la donnée de descente sur $X'$.} Il en est
certainement ainsi (quels que soient $X'/S'$ et la donnée de descente sur $X'$)
si le morphisme $\alpha:S'\longrightarrow S$ est \emph{radiciel} (i.e.
injectif, et à extensions résiduelles qui

\src{318}
sont radicielles). On peut montrer aussi qu'il en est encore ainsi si
$\alpha:S'\longrightarrow S$ est \emph{fini}, et toute partie finie de $X'$,
contenue dans une fibre de $X'$ sur $S$, est contenue dans un ouvert de $X'$
affine sur $S$ (c'est le \emph{critère de Weil}). Il en est en particulier
ainsi si $X'/S'$ est quasi-projectif, et dans ce cas on peut montrer que le
préschéma « descendu » $X/S$ est aussi quasi-projectif (et projectif si $X'/S'$
l'est). En résumé :

\result{THÉORÈME}{3} Soit $\alpha:S'\longrightarrow S$ un morphisme de
préschémas, fidèlement plat et quasi-compact. Si $\alpha$ est \emph{radiciel},
c'est un \emph{morphisme de descente strict}. Si $\alpha$ est fini, c'est un
morphisme de descente strict relativement à la catégorie fibrée des préschémas
quasi-projectifs (ou projectifs) sur des préschémas.

\result{REMARQUES}{} J'ignore si, dans la deuxième assertion ci-dessus,
l'hypothèse que $\alpha$ soit un morphisme \emph{fini} est bien nécessaire ; on
vérifie en tout cas de façon « formelle » qu'on peut la remplacer par
l'hypothèse suivante, plus faible en apparence : \emph{pour tout point de $S$
existe un voisinage ouvert $U$, un $U'$ fini et fidèlement plat sur $U$, et un
$S$-morphisme de $U'$ dans $S'$.} Un cas type qui ne rentre pas dans le
précédent est celui où $S=\operatorname{Spec}(A)$,
$S'=\operatorname{Spec}(\widehat A)$, où $A$ est un anneau local noethérien et
$\widehat A$ son complété ; ou encore celui où $S'$ est quasi-fini sur $S$
(i.e. localement isomorphe à un ouvert d'un $S$-schéma fini) et non fini. Dans
ces deux cas, le conférencier ignore aussi la réponse à la question suivante :
soit $X$ un $S$-schéma tel que $X'=X\times_SS'$ soit projectif sur $S'$, est-il
vrai que $X$ est projectif sur $S$ ?\footnote{L'erratum de 1962 répond
négativement : un morphisme $S'\longrightarrow S$ quasi-fini, étale et
surjectif, ou un morphisme
$\operatorname{Spec}(\widehat A)\longrightarrow\operatorname{Spec}(A)$, n'est
pas toujours un morphisme de descente strict, même si $A$ est l'anneau local
d'une courbe algébrique sur un corps algébriquement clos $k$ et
$S=\operatorname{Spec}(A)$. On peut ainsi trouver un morphisme propre et simple
$f:X\longrightarrow S$, faisant de $X$ un fibré principal de base $S$ sous une
courbe elliptique $E$, tel que $f':X'\longrightarrow S'$ soit projectif, mais
$f$ non projectif. C'est donc aussi un exemple de fibré principal homogène non
isotrivial sous un schéma abélien.}

\subsection*{2. Application à la descente de certaines propriétés de morphismes}
\addcontentsline{toc}{subsection}{2. Descente de propriétés de morphismes}

Soit $P$ une classe de morphismes de préschémas. Soit
$\alpha:S'\longrightarrow S$ un morphisme de préschémas, et soit
$f:X\longrightarrow Y$ un morphisme de $S$-préschémas,
$f':X'\longrightarrow Y'$ l'image réciproque de $f$ par $\alpha$. On peut se
demander alors si la relation « $f'\in P$ » implique « $f\in P$ ». Il apparaît
que la réponse est affirmative dans beaucoup de cas importants lorsque l'on
suppose que $\alpha$ est \emph{fidèlement plat et quasi-compact} (cette
dernière hypothèse étant parfois surabondante). Cela se voit directement sans
difficulté si $P$ est la classe des morphismes surjectifs, resp. radiciels (ces
cas résultant de la surjectivité de $\alpha$), resp. plats, resp. fidèlement
plats, resp. simples (ces cas résultant de la fidèle platitude de $\alpha$),
resp. de type fini. Utilisant les théorèmes 1, 2 et le lemme 1.2, on voit aussi
qu'il en est de même si $P$ est une des classes suivantes : isomorphismes,
immersions ouvertes, immersions fermées, immersions (si $f$ est de type fini et
$Y$ localement noethérien), morphismes affines, morphismes finis, morphismes
quasi-finis, morphismes ouverts, morphismes fermés, homéomorphismes, morphismes
séparés,

\src{319}
morphismes propres. Le seul cas important non élucidé est celui des morphismes
projectifs et quasi-projectifs, déjà signalé dans la remarque du paragraphe 1.

\subsection*{3. Descente par morphismes finis fidèlement plats}
\addcontentsline{toc}{subsection}{3. Descente par morphismes finis fidèlement plats}

Soit $\alpha:S'\longrightarrow S$ un morphisme fini, correspondant à un
faisceau d'algèbres $\mathcal A'$ sur $S$ qui est localement libre de type fini
en tant que faisceau de modules, et partout $\ne0$ ; alors $\alpha$ est un
morphisme fidèlement plat et quasi-compact, auquel on peut donc appliquer les
résultats précédents. La donnée d'un faisceau quasi-cohérent $F'$ sur $S'$
équivaut à la donnée du faisceau quasi-cohérent $\alpha_*(F')$ sur $S$, muni de
sa structure de $\mathcal A'$-module (notant que
$\mathcal A'=\alpha_*(\mathcal O_{S'})$). Pour simplifier, ce faisceau sur $S$
sera également noté $F'$. Les deux images inverses $p_i^*(F')$ de $F'$ sur
$S'\times_SS'$ correspondent de même aux faisceaux quasi-cohérents de
$(\mathcal A'\otimes_{\mathcal O_S}\mathcal A')$-modules
\[
  F'\otimes_{\mathcal O_S}\mathcal A'
  \quad\text{et}\quad
  \mathcal A'\otimes_{\mathcal O_S}F'.
\]
La donnée d'un $(S'\times_SS')$-homomorphisme du premier dans le second équivaut
à la donnée d'un homomorphisme de
$(\mathcal A'\otimes_{\mathcal O_S}\mathcal A')$-modules et, compte tenu que
$\mathcal A'$ est localement libre, cela équivaut à la donnée d'un
homomorphisme de $(\mathcal A'\otimes\mathcal A')$-modules :
\[
  \mathcal U=
  \operatorname{Hom}_{\mathcal O_S}(\mathcal A',\mathcal A')
  =\mathcal A'\otimes\check{\mathcal A}'
  \longrightarrow
  \operatorname{Hom}_{\mathcal O_S}(F',F'),
\]
i.e. à la donnée, pour toute section $\xi$ de $\mathcal U$ sur un ouvert $V$,
d'un homomorphisme de $\mathcal O_S$-modules
$T_\xi:F'|_V\longrightarrow F'|_V$, satisfaisant les conditions
\begin{equation}
  T_{f\xi}(x)=fT_\xi(x),
  \qquad T_{\xi f}(x)=T_\xi(fx),
  \tag{3.1}
\end{equation}
où $f,x$ sont respectivement des sections de $\mathcal A',F'$ sur un ouvert de
$S$ contenu dans $V$. Les conditions (i) et (ii) d'une donnée de descente
(A, paragraphe 1, (c)) s'écrivent alors respectivement :
\begin{equation}
  T_{1_{\mathcal U}}x=x,
  \qquad\text{i.e.}\qquad
  T_{1_{\mathcal U}}=\operatorname{id}_{F'},
  \tag{3.2}
\end{equation}
\begin{equation}
  T_{\xi\eta}=T_\xi T_\eta.
  \tag{3.3}
\end{equation}
En d'autres termes, \emph{une donnée de descente sur $F'$ équivaut à une
représentation du faisceau de $\mathcal O_S$-algèbres
$\mathcal U=\operatorname{Hom}_{\mathcal O_S}(\mathcal A',\mathcal A')$ dans
le faisceau de $\mathcal O_S$-algèbres
$\operatorname{Hom}_{\mathcal O_S}(F',F')$, satisfaisant les deux conditions de
linéarité (3.1).} Si on a un accouplement de faisceaux quasi-cohérents sur
$S'$ :
\[
  F'_1\times F'_2\longrightarrow F'_3
\]
(qu'on peut interpréter comme un accouplement de faisceaux de
$\mathcal A'$-modules sur $S$),

\src{320}
et des données de recollement sur les $F'_i$, définies par des homomorphismes
\[
  T_i:\mathcal U\longrightarrow
  \operatorname{Hom}_{\mathcal O_S}(F'_i,F'_i)
  \qquad (i=1,2,3),
\]
alors ces données sont compatibles avec l'accouplement donné, dans le sens
évident du terme, si et seulement si la condition suivante est vérifiée.

Pour toute section $\xi$ de $\mathcal U$ sur un ouvert, désignant par
\[
  \Delta\xi=\sum_i\xi'_i\otimes_{\mathcal A'}\xi''_i
\]
la section de $\mathcal U\otimes_{\mathcal A'}\mathcal U$ ($\mathcal U$ étant
considéré comme $\mathcal A'$-module pour sa structure à gauche) définie par la
formule
\[
  \xi\mathbin{\cdot}(fg)=\sum_i\xi'_i(f)\xi''_i(g)
\]
(où $f$ et $g$ sont deux sections de $\mathcal A'$ sur un ouvert plus petit),
on a la formule
\begin{equation}
  T^{(3)}_\xi(x\mathbin{\cdot}y)
  =\sum_iT^{(1)}_{\xi'_i}x\mathbin{\cdot}T^{(2)}_{\xi''_i}y
  \tag{3.4}
\end{equation}
pour tout couple de sections $x,y$ de $F'_1,F'_2$ respectivement sur un ouvert
plus petit.\footnote{Le texte imprimé porte « de $\mathcal A'$ ». Les domaines
de $T^{(1)}$ et $T^{(2)}$, ainsi que l'accouplement
$F'_1\times F'_2\longrightarrow F'_3$, imposent que $x$ et $y$ soient des
sections de $F'_1$ et $F'_2$ respectivement.} (On pourra exprimer cette
propriété en disant que les homomorphismes $T^{(i)}$ sont compatibles avec
l'application diagonale de $\mathcal U$, relativement à l'accouplement donné.)
En particulier, les formules (3.1) à (3.4) nous permettent d'interpréter en
termes de représentations d'algèbres à applications diagonales les données de
descente sur un faisceau quasi-cohérent d'algèbres sur $S'$, donc aussi (se
restreignant aux algèbres commutatives) les données de descente sur un
$S'$-schéma affine.

De là, on passe à une interprétation analogue des données de descente sur un
$S'$-pré-schéma $X'$ quelconque : la donnée d'un tel $X'$ équivaut à la donnée
d'un pré-schéma $X'$ sur $S$, muni d'un homomorphisme de $\mathcal O_S$-algèbres
\[
  \mathcal A'\longrightarrow\mathcal O_{X'},
\]
et une donnée de descente sur $X'$ équivaut à la donnée d'un homomorphisme de
faisceaux (compatible avec le morphisme $h:X'\longrightarrow S'$)
\[
  \mathcal U\longrightarrow
  \operatorname{Hom}_{h^{-1}(\mathcal O_S)}
  (\mathcal O_{X'},\mathcal O_{X'}),
\]
satisfaisant les conditions analogues aux conditions (3.1) à (3.4) ci-dessus.

\textsc{Exemple 1 (« Weil »).} -- Supposons $S'/S$ un revêtement étale
galoisien de groupe de Galois $\Gamma$ (cf. A, paragraphes 3 et 4.d). Alors une
donnée de descente sur un faisceau quasi-cohérent $F'$ sur $S'$ (resp. sur un
$S'$-pré-schéma $X'$) équivaut à la donnée d'une représentation de $\Gamma$ par
automorphismes de $(S',F')$ (resp. de $(S',X')$) compatible avec les opérations
de $\Gamma$ sur $S'$. Ce résultat

\src{321}
est « formel », i.e. se démontre en termes de catégories, mais du point de vue
de ce numéro résulte aussi de la structure explicite de $\mathcal U$ (muni de
sa structure d'anneau, l'homomorphisme d'anneaux
$\mathcal A'\longrightarrow\mathcal U$ et l'application diagonale),
complètement connue grâce au résultat suivant : $\mathcal U$ admet, en tant que
$\mathcal A'$-module à gauche, une base formée des sections de $\mathcal U$ qui
correspondent aux éléments de $\Gamma$.

\textsc{Exemple 2 (« Cartier »).} -- Soit $p$ un nombre premier, supposons
$p\mathcal O_S=0$ (i.e. $\mathcal O_S$ est de caractéristique $p$),
$(\mathcal A')^p\subset\mathcal O_S=\mathcal A$ (i.e. $S'/S$ est radiciel de
hauteur 1) et que le faisceau d'algèbres $\mathcal A'$ sur $\mathcal A$ admet
localement une $p$-base, i.e. une famille $(x_i)$ de sections telle que
$\mathcal A'$ soit engendré comme algèbre par les $x_i$, soumis aux seules
conditions $x_i^p=0$. Nous supposons l'ensemble des $i$ fini, de cardinal $n$.
Soit $\mathfrak g$ le faisceau des $\mathcal A$-dérivations de $\mathcal A'$ ;
c'est un faisceau de $\mathcal A'$-modules localement libre de rang $n$, de plus
c'est un faisceau de $p$-algèbres de Lie sur $\mathcal A$ (mais non sur
$\mathcal A'$) satisfaisant la condition
\begin{equation}
  [X,fY]=X(f)Y+f[X,Y].
  \tag{3.5}
\end{equation}

\textsc{Lemme.} -- $\mathcal U=
\operatorname{Hom}_{\mathcal O_S}(\mathcal A',\mathcal A')$ est engendré, en
tant que $\mathcal O_S$-algèbre muni d'un homomorphisme d'algèbres
$\mathcal A'\longrightarrow\mathcal U$, par le sous-$\mathcal A'$-module à
gauche $\mathfrak g$, avec les relations supplémentaires
\begin{equation}
  \left\{
  \begin{aligned}
    Xf-fX&=X(f),\\
    XY-YX&=[X,Y],\\
    X^p&=X^{(p)}.
  \end{aligned}
  \right.
  \tag{3.6}
\end{equation}

Il résulte du lemme précédent qu'une donnée de descente sur le faisceau
quasi-cohérent $F'$ sur $S'$ équivaut à la donnée, pour tout $X\in\mathfrak g$,
d'un $\mathcal O_S$-endomorphisme $\overline X$ de $F'$, satisfaisant les
conditions
\begin{align}
  \overline{fX}&=f\overline X, \tag{3.7}\\
  \overline X(fx)&=X(f)x+f\overline X(x), \tag{3.8}\\
  \overline{[X,Y]}&=[\overline X,\overline Y], \tag{3.9}\\
  \overline{X^{(p)}}&=\overline X^p. \tag{3.10}
\end{align}
(C'est ce qu'on pourrait appeler une connexion linéaire sur $F'$, sans
courbure et compatible avec la puissance $p$-ième.) On explicite de même la
notion de donnée de descente sur un $S'$-pré-schéma $X'$ ; la relation (3.4) est
remplacée ici par la condition que les $\overline X$ sont des dérivations de
$\mathcal O_{X'}$. Comme le morphisme $S'\longrightarrow S$ est radiciel, le
théorème 3 garantit que toute telle donnée de descente est

\src{322}
effective, donc définit un $S$-pré-schéma $X$.

On notera qu'on n'a pas eu à faire d'hypothèse de platitude, de non singularité
ou de finitude quelconque sur $F'$ resp. $X'$.

\subsection*{4. Application à des critères de rationalité}
\addcontentsline{toc}{subsection}{4. Application à des critères de rationalité}

Soit $X$ un $S$-pré-schéma tel que l'image directe de $\mathcal O_X$ sur $S$
soit $\mathcal O_S$ ; cette propriété restera vraie alors par toute extension
plate $S'\longrightarrow S$ de la base $S$. Si $F$ est un faisceau inversible
(i.e. localement libre de rang 1) sur $X$, les automorphismes de $F$,
s'identifiant aux sections inversibles de $\mathcal O_X$, correspondent
biunivoquement aux sections inversibles de $\mathcal O_S$. Soit alors $s$ une
section de $X$ au-dessus de $S$ ; nous appellerons de façon imagée, section de
$F$ au-dessus de $s$, une section du faisceau inversible $s^*(F)$ sur $S$. Il
résulte de ce qui précède que si $F_i$ ($i=1,2$) sont deux faisceaux inversibles
sur $X$, munis chacun d'une section au-dessus de $s$, et si $F_1$ et $F_2$ sont
isomorphes, il existe un isomorphisme et un seul de $F_1$ sur $F_2$ compatible
avec les sections en question (i.e. transformant la première en la seconde).
D'ailleurs, et indépendamment de la section $s$, convenons de regarder comme
équivalents deux faisceaux inversibles $F_1$ et $F_2$ sur $X$ tels que tout
point de $S$ ait un voisinage ouvert $U$ tel que les restrictions de $F_1$ et
$F_2$ à $X|U$ soient isomorphes. Alors tout faisceau inversible $F$ sur $X$ est
équivalent à un faisceau inversible $F_1$ muni d'une section marquée au-dessus
de $s$ (on prend $F_1=F\,s^*(F)^{-1}$), et $F_1$ est déterminé à un
isomorphisme près. En d'autres termes, la classification des faisceaux
inversibles sur $X$ à une équivalence près est la même que la classification à
un isomorphisme près des faisceaux inversibles munis d'une section marquée.

Comme ces propriétés restent vraies par extension plate
$\alpha:S'\longrightarrow S$ de la base (en remplaçant la section $s$ par son
image inverse $s'$ par $\alpha$), on en conclut, compte tenu du théorème 1 :

\emph{Le pré-schéma $X/S$ étant comme ci-dessus et admettant une section $s$,
soit $\alpha:S'\longrightarrow S$ un morphisme fidèlement plat et
quasi-compact ; soit $F'$ un faisceau inversible sur $X'=X\times_SS'$. Pour que
$F'$ soit équivalent à l'image inverse sur $X'$ d'un faisceau inversible $F$
sur $X$, il faut et il suffit que ses images inverses $p_1^*(F')$ et $p_2^*(F')$
sur $X'\times_XX'=X\times_S(S'\times_SS')$ soient équivalentes. S'il en est
ainsi, $F$ est déterminé à une équivalence près.} (On dira alors que $F'$ est
rationnel sur $S$.)

S'inspirant de ce principe dans le cas où
$\alpha:S'\longrightarrow S$ est comme dans l'exemple 1 ou l'exemple 2 du
numéro précédent, on trouve les critères de rationalité de Weil ou de Cartier.
(On notera que ces auteurs se bornent au cas où $S$ et $S'$

\src{323}
sont des spectres de corps ; a fortiori, $S$ est alors le spectre d'un anneau
local, et la relation d'équivalence introduite ci-dessus n'est autre que la
relation d'isomorphie.) Dans le premier cas, $F'$ est rationnel sur $S$ si et
seulement si ses transformés par $\Gamma$ sont équivalents à $F'$. Pour
exprimer le critère de rationalité dans le deuxième cas, on considère, de façon
générale, le morphisme diagonal
\[
  X'\longrightarrow X''=X'\times_XX'
\]
de $X'/X$, le faisceau d'idéaux $\mathcal I$ correspondant sur
$X'\times_XX'$ et le faisceau $\mathcal I/\mathcal I^2$, qui s'identifie à son
image réciproque $\Omega^1_{X'/X}$ sur $X'$. Comme les restrictions des
$F''_i=p_i^*(F')$ ($i=1,2$) à la diagonale sont isomorphes (car isomorphes à
$F'$), i.e. $F''_1(F''_2)^{-1}=F''$ a une restriction à la diagonale qui est
triviale, il s'ensuit que la restriction de $F''$ à
$(X'',\mathcal O_{X''}/\mathcal I^2)$ est donnée, à un isomorphisme près, par
un élément bien déterminé $\xi$ de
\[
  H^1(X'',\mathcal I/\mathcal I^2)
  =H^1(X',\Omega^1_{X'/X}).
\]
D'ailleurs, en l'occurrence, on a
\[
  \Omega^1_{X'/X}
  =\Omega^1_{S'/S}\otimes_{\mathcal O_S}\mathcal O_X,
\]
et par suite, si $\Omega^1_{S'/S}$ est localement libre sur $S$ (comme dans le
cas de Cartier), $\xi$ définit une section de
\[
  R^1f'_*(\mathcal O_{X'})\otimes\Omega^1_{S'/S}
\]
sur $S'$ (appelée classe de Atiyah--Cartier du faisceau inversible $F'$ sur
$X'/S$), dont l'annulation est nécessaire et suffisante pour que les images
réciproques de $F'$ par les deux projections de
\[
  (X'',\mathcal O_{X''}/\mathcal I^2)
  =X\times_S(S'',\mathcal O_{S''}/\mathcal J^2)
\]
sur $X'$ soient équivalentes (où $\mathcal J$ est le faisceau d'idéaux sur
$S''=S'\times_SS'$ défini par le morphisme diagonal
$S'\longrightarrow S'\times_SS'$). Cette annulation est donc trivialement
nécessaire pour que les images réciproques de $F'$ sur
$X''=X\times_SS''$ lui-même soient équivalentes, donc aussi pour que $F'$ soit
équivalent à l'image inverse d'un faisceau inversible $F$ sur $X$. D'ailleurs,
la classe de Atiyah--Cartier peut aussi s'interpréter comme l'obstruction à
l'existence, localement au-dessus de $S'$, d'une connexion de $F'$ relativement
aux dérivations de $X'/X$, une telle connexion étant de plus déterminée quand
on connaît les dérivations de $F'$ correspondant aux prolongements naturels à
$X'$ des dérivations de $S'/S$. De ceci, et des développements du numéro
précédent, on conclut facilement que dans le cas de l'exemple 2 dudit numéro,
et lorsque $X/S$ admet une section, l'annulation de la classe de
Atiyah--Cartier est aussi suffisante pour que $F'$ soit rationnelle sur $S$.

\subsection*{5. Application à la restriction du schéma de base dans un schéma abélien}
\addcontentsline{toc}{subsection}{5. Restriction du schéma de base dans un schéma abélien}

Soit $S$ un pré-schéma. On appelle schéma abélien sur $S$ un schéma $X$ simple
et propre

\src{324}
sur $S$ dont les fibres en les points $x\in S$ sont des schémas de variétés
abéliennes sur les $\kappa(x)$. Supposons $S$ noethérien et régulier (i.e. ses
anneaux locaux réguliers), alors on peut montrer en utilisant le théorème de
connexion de Murre [4] (du moins dans le cas « d'égales caractéristiques », où
le théorème cité est actuellement démontré) que toute section rationnelle de
$X$ sur $S$ est partout définie (i.e. est une section) (ce qui généralise un
théorème classique de Weil). Il en résulte, plus généralement, que si $X'$ est
un schéma simple sur $S$, alors toute $S$-application rationnelle de $X'$ dans
$X$ est partout définie. Il en résulte ceci, qui généralise un résultat de
Chow--Iang :\footnote{Le texte imprimé porte « Igusa--Lang » ; l'erratum de
1962 prescrit « Chow--Iang ».} $S$ étant noethérien et régulier et $K$ désignant
son anneau des fonctions rationnelles (composé direct de corps), soit $X$ un
schéma abélien au-dessus de $K$ ; si $X$ est isomorphe à un $K$-schéma de la
forme $X_0\times_S\operatorname{Spec}(K)$, où $X_0$ est un schéma abélien sur
$S$, alors $X_0$ est déterminé à un isomorphisme unique près.

Utilisant le résultat d'unicité précédent, on voit que la question de
restriction de la base dans $X$ est locale sur $S$ (et par suite, qu'il suffit
de savoir faire la restriction aux $\operatorname{Spec}(\mathcal O_x)$, avec
$x\in S$). On voit de la même façon que si
$\alpha:S'\longrightarrow S$ est un morphisme simple de type fini, si $K'$ est
l'anneau des fonctions rationnelles de $S'$, et si $X\otimes_KK'$ est de la
forme $X'_0\times_{S'}\operatorname{Spec}(K')$, alors $X'_0$ est muni d'une
donnée de descente canonique relativement à $\alpha$. Compte tenu du théorème
3, on en conclut :

\result{PROPOSITION}{5.1} Soit $S$ un pré-schéma noethérien et régulier,
irréductible, de corps de fonctions rationnelles $K$, soit $K'$ une extension
finie de $K$, non ramifiée sur $S$, $S'$ le normalisé de $S$ dans $K'$ (qui est
donc un revêtement étale de $S$), $X$ un schéma abélien sur $K$ tel que
$X\otimes_KK'$ soit de la forme
$X'_0\times_{S'}\operatorname{Spec}(K')$, où $X'_0$ est un schéma abélien
projectif sur $S'$. Alors $X$ est de la forme
$X_0\times_S\operatorname{Spec}(K)$, où $X_0$ est un schéma abélien projectif
sur $S$.

\result{REMARQUES}{} Le conférencier ignore si on peut remplacer l'hypothèse
que $S'\longrightarrow S$ est un revêtement étale surjectif (permettant
d'utiliser le théorème 3) par l'hypothèse que c'est un morphisme de type fini
simple et surjectif (même si on suppose que c'est un étalement), ou si la
proposition reste valable sans supposer $X'_0$ projectif sur $S'$ (condition qui
est peut-être remplie automatiquement).

\subsection*{6. Application à des critères de locale trivialité et isotrivialité}
\addcontentsline{toc}{subsection}{6. Critères de locale trivialité et isotrivialité}

Soient $S$ un pré-schéma, $G$ un « pré-schéma en groupes » au-dessus de $S$,
$P$ un pré-schéma sur $S$ sur lequel « $G$ opère » (à droite). On dit que $P$
est formellement principal homogène sous $G$ si le morphisme bien connu
\[
  G\times_SP\longrightarrow P\times_SP
\]

\src{325}
déduit des opérations de $G$ sur $P$, est un isomorphisme. Nous supposerons
dorénavant $G$ plat sur $S$ (donc fidèlement plat sur $S$), et nous réserverons
le nom de fibré principal homogène sous $G$ à un fibré formellement principal
homogène $P$ qui est fidèlement plat et quasi-compact sur $S$. Il est immédiat
qu'il revient au même de dire que l'on peut trouver une extension fidèlement
plate et quasi-compacte $S'\longrightarrow S$ de la base $S$, telle que le
fibré formellement principal homogène
\[
  P'=P\times_SS'
\]
sous $G'=G\times_SS'$ soit trivial, i.e. isomorphe à $G'$ (i.e. admette une
section) ; on pourra prendre en particulier $S'=P$. Noter d'ailleurs que si
$S$ est localement noethérien, alors l'hypothèse de fidèle platitude sur $P$
équivaut à l'hypothèse que
\[
  \widehat P_s=P\times_S\operatorname{Spec}(\widehat{\mathcal O}_s)
\]
est fidèlement plat sur $\widehat{\mathcal O}_s$ pour tout $s\in S$ (où
$\widehat{\mathcal O}_s$ désigne le complété de l'anneau local
$\mathcal O_s$), comme il résulte du fait que $\widehat{\mathcal O}_s$ est
fidèlement plat sur $\mathcal O_s$. D'ailleurs, si $P$ est de type fini sur
$S$ localement noethérien, l'ensemble des points $s$ satisfaisant à la
condition ci-dessus est constructible, donc si $S$ est un « pré-schéma de
Jacobson » (par exemple un schéma de type fini sur un corps, ou le spectre d'un anneau de
Jacobson plus généralement), il suffit de vérifier la condition en question
pour les points fermés de $S$. Cela nous ramène au cas où la base est le
spectre d'un anneau local complet $A$. Lorsque $S=\operatorname{Spec}(A)$
($A$ anneau local noethérien complet) et que $P$ est de type fini sur $S$, la
fidèle platitude de $P/S$ équivaut aussi à l'existence d'un $S'$ fini et plat
sur $S$ tel que $P'$ soit trivial, et si de plus $G$ est simple sur $S$, on
peut supposer $S'$ étalé sur $S$. Par suite, si de plus le corps résiduel de
$A$ est algébriquement clos (« cas géométrique »), $P$ est fidèlement plat sur
$A$ si et seulement s'il est trivial. Donc, si $S$ est un pré-schéma
algébrique sur un corps algébriquement clos, et $G$ simple de type fini sur
$S$, on voit que la condition de fidèle platitude sur $S$ équivaut à la
condition de trivialité analytique (SLF) de Serre ([6], p. 1--12).

On peut introduire d'autres types plus forts de conditions sur $P$, ayant la
nature d'une « trivialité locale ». Nous dirons en particulier que $P$ est
isotrivial (resp. strictement isotrivial) si pour tout $s\in S$, il existe un
voisinage ouvert $U$ de $s$, et un morphisme fini et fidèlement plat (resp. un
revêtement étale surjectif) $U'\longrightarrow U$ tel que
$P'=P\times_SU'$ soit trivial. (Nous nous écartons de la terminologie de Serre
[6], qui appelle localement isotrivial ce que nous appelons strictement
isotrivial.) La stricte isotrivialité est surtout utile si $G$ est simple sur
$S$, mais est une notion inadéquate par contre dans les autres cas.

Si $G$ est affine sur $S$, tout fibré principal homogène $P$ sous $G$ est
affine d'après le paragraphe 2, d'où la possibilité, grâce au théorème 2, de
« descendre »

\src{326}
de tels fibrés par des morphismes fidèlement plats et quasi-compacts. Prenons,
en particulier, $G=\operatorname{Gl}(n)_S$, défini par la condition que le
foncteur de $S$-pré-schémas
\[
  S'\longmapsto\operatorname{Hom}_S(S',G)
\]
(à valeurs dans la catégorie des groupes) s'identifie au foncteur
\[
  \operatorname{Gl}(n)(S')
  =\operatorname{Gl}\bigl(n,H^0(S',\mathcal O_{S'})\bigr)
\]
du A, paragraphe 4.e. Utilisant le fait (i) que tout fibré principal homogène
sous $G$ (resp. tout faisceau localement libre de rang $n$ sur $S$) devient
isomorphe à l'objet « trivial » $G$ (resp. $\mathcal O_S^n$) par une extension
fidèlement plate et quasi-compacte convenable de $S$ ; (ii) que l'on peut
descendre les objets du type envisagé (fibrés principaux homogènes sous $G$,
resp. faisceaux localement libres de rang $n$) par de tels morphismes, et enfin
(iii) que le groupe des automorphismes du fibré trivial sur un $S'/S$ est
fonctoriellement isomorphe au groupe des automorphismes du faisceau localement
libre de rang $n$ trivial sur $S'$, on conclut « formellement » qu'il « revient
au même » de se donner sur $S$ (ou sur un $S'/S$) un fibré principal homogène
de groupe $G$, ou de s'y donner un faisceau localement libre de rang $n$. (De
façon plus précise, on a une équivalence de catégories fibrées.) On en conclut
en particulier :

\result{PROPOSITION}{6.1} Tout fibré principal homogène de groupe
$\operatorname{Gl}(n)_S$ est localement trivial.

Par des arguments connus, on en conclut le même résultat pour des groupes
structuraux tels que $\operatorname{Sl}(n)_S$, $\operatorname{Sp}(n)_S$ et des
produits de tels groupes. On en conclut aussi que, si $F$ est un sous-groupe
fermé de $G=\operatorname{Gl}(n)_S$, plat sur $S$, tel que le quotient $G/F$
existe, et que $G$ soit un fibré principal homogène isotrivial (resp.
strictement isotrivial) sur $G/F$, de groupe $F\times_S(G/F)$, alors tout fibré
principal homogène de groupe $F$ est isotrivial (resp. strictement isotrivial).
Cela s'applique à tous les « groupes linéaires » sur $S$ qui ont été utilisés
jusqu'à présent, et en particulier au cas où $G=S\times_k\Gamma$, $S$ étant un
pré-schéma sur le corps $k$, et $\Gamma$ un groupe linéaire au sens classique,
et en particulier simple, sur $k$. Cela résout donc, pour de tels groupes, une
question de Serre (loc. cit.).

Signalons aussi que, pour la plupart des groupes (linéaires ou non) simples sur
$S$ connus, et en tout cas ceux de la forme $S\times_k\Gamma$ comme ci-dessus,
on peut montrer que tout fibré principal homogène isotrivial est strictement
isotrivial, ce qui résout en particulier une autre question de Serre (loc.
cit., p. 1--14), compte tenu qu'un fibré principal homogène obtenu par une
descente à la Cartier (cf. paragraphe 3, exemple 2) est manifestement
isotrivial.

\result{REMARQUE}{} Une des difficultés essentielles dans ces questions (mise
à part la question de l'existence des schémas quotients $G/F$) est le manque de
critères d'effectivité pour une donnée de descente par un morphisme fidèlement
plat non fini.

\src{327}
\begin{center}
  \textbf{BIBLIOGRAPHIE}
\end{center}

\begin{description}
  \item[{[1]}] Dieudonné (J.) et Grothendieck (Alexander). --- \emph{Éléments
    de géométrie algébrique}, à paraître dans les Publications mathématiques de
    l'Institut des Hautes Études Scientifiques.
  \item[{[2]}] Grauert (H.) und Remmert (R.). --- \emph{Komplexe Räume}, Math.
    Annalen, t. 136, 1958, p. 245--318.
  \item[{[3]}] Grothendieck (Alexander). --- \emph{Géométrie formelle et
    géométrie algébrique}. Séminaire Bourbaki, t. 11, 1958/59, no 182.
  \item[{[4]}] Murre (J. P.). --- \emph{On a connectedness theorem for a
    birational transformation at a simple point}, Amer. J. Math., t. 80, 1958,
    p. 3--15.
  \item[{[5]}] Serre (Jean-Pierre). --- \emph{Géométrie algébrique et géométrie
    analytique}, Ann. Inst. Fourier Grenoble, t. 6, 1955--56, p. 1--42.
  \item[{[6]}] Serre (Jean-Pierre). --- \emph{Espaces fibrés algébriques},
    Séminaire Chevalley, t. 2, 1958 : Anneaux de Chow et applications, no 1.
\end{description}

\par\smallskip
\begingroup\footnotesize
\endgroup
